% answers_short.tex —— 2008 年·普通高等学校招生考试·新课标卷（琼、宁卷）·理科 · 简短答案速查
% 由 tools/make_short_answers.py 自 answers.tex 的 \gksolhead/\gkanswer 提取，
% 并与 2008/2008_数学真题_新课标理_答案要点.md 逐题比对后生成，勿手改（改源 MD 后重跑）。

\begin{gknotice}
\textbf{使用说明}：本节只给最终答案便于快速核对，完整推导见解析版。答案为编写者独立推导，\textbf{非官方标准答案}。
\end{gknotice}

\gkpart{选择题}
\begin{center}\begin{tabular}{ccccccccccccc}
\toprule
题号 & 1 & 2 & 3 & 4 & 5 & 6 & 7 & 8 & 9 & 10 & 11 & 12 \\
\midrule
答案 & B & B & D & C & A & B & C & D & A & D & A & C \\
\bottomrule
\end{tabular}\end{center}

\gkpart{填空题}
\begin{itemize}[leftmargin=2.6em]
  \item[\textbf{13.}] $3$。
  \item[\textbf{14.}] $\dfrac{32}{15}$。
  \item[\textbf{15.}] $\dfrac{4\pi}3$。
  \item[\textbf{16.}] （开放性作答，以下两条均可）
\end{itemize}

\gkpart{解答题}
\begin{itemize}[leftmargin=2.6em]
  \item[\textbf{17.}] （1）$a_n=5-2n$；（2）$S_n$ 的最大值为 $4$。
  \item[\textbf{18.}] （1）$DP$ 与 $CC'$ 所成角为 $45^\circ$；（2）$DP$ 与平面 $AA'D'D$ 所成角为 $30^\circ$。
  \item[\textbf{19.}] （1）$DY_1=4$，$DY_2=12$；（2）$x=75$ 时 $f(x)$ 取最小值 $3$。
  \item[\textbf{20.}] （1）$\dfrac{x^2}4+\dfrac{y^2}3=1$；（2）$y=\sqrt6x+2\sqrt3$ 或 $y=\sqrt6x-2\sqrt3$。
  \item[\textbf{21.}] （1）$f(x)=x+\dfrac1{x-1}$；（2）对称中心 $(1,1)$；（3）三角形面积为定值 $2$。
\end{itemize}

\gkpart{选做题}
\begin{itemize}[leftmargin=2.6em]
  \item[\textbf{22.}] （1）（2）证明见解析。
  \item[\textbf{23.}] （1）$C_1$ 是单位圆 $x^2+y^2=1$，$C_2$ 是直线 $y=x+\sqrt2$，二者有 $1$ 个公共点；（2）$C_1':x^2+4y^2=1$，$C_2':y=\dfrac{x+\sqrt2}2$，公共点个数相同，仍为 $1$。
  \item[\textbf{24.}] （1）图象见解析；（2）$\{x\mid x<5\}$。
\end{itemize}
